\section{Irreducibility and symmetric monodromy}\label{sec:monodromy}

Let $k$ be any field of characteristic zero and let $H\in k[W]$ have degree
$n\ge2$; admissibility at $0$ and $1$ is not required in this section.  Put
\begin{equation}\label{eq:universal-pencil}
 \mathscr E(W)=H(W)-sW+t\in k[s,t,W].
\end{equation}

\begin{proposition}[geometric irreducibility]\label{prop:pencil-irreducible}
The polynomial $\mathscr E$ is irreducible over $\overline{k}(s,t)$ and over
$k(s,t)$.
\end{proposition}

\begin{proof}
Work first in $\overline{k}[s,W,t]$.  Since $\mathscr E$ has degree one in
$t$, in any factorization one factor, say $A$, has $t$-degree zero.  Write the
other factor as $B_0+B_1t$.  Comparison of the coefficient of $t$ gives
$AB_1=1$ in $\overline{k}[s,W]$, so $A$ and $B_1$ are units.  Thus the
factorization is trivial.  The leading coefficient in $W$ is the nonzero
constant leading coefficient of $H$, so Gauss's lemma gives irreducibility
over $\overline{k}(s,t)$.  The assertion over $k(s,t)$ follows.
\end{proof}

\begin{proposition}[discriminant normalization]\label{prop:universal-disc}
The repeated-root discriminant of \eqref{eq:universal-pencil} is an
irreducible curve whose normalization is birationally parametrized by
\begin{equation}\label{eq:disc-param}
 r\longmapsto\left(s,t\right)=\left(H'(r),rH'(r)-H(r)\right).
\end{equation}
Its generic point represents a polynomial with exactly one double root and
all other roots simple.
\end{proposition}

\begin{proof}
A repeated root $r$ satisfies $\mathscr E(r)=\mathscr E'(r)=0$, which is
equivalent to \eqref{eq:disc-param}.  The two coordinate functions have pole
orders $n-1$ and $n$ at $r=\infty$.  The generic degree of the induced map
from $\P^1$ to its image divides both pole orders, hence is one.  The image is
therefore irreducible and the map is birational.  The morphism from $\A^1_r$
is finite because $r$ is integral over $k[H'(r)]$: after division by the
leading coefficient, $H'(T)-H'(r)$ is a monic equation for $r$ over that
subring.  Since $\A^1$ is normal, it is the normalization of the affine
discriminant.  Since $H''$ is not the zero
polynomial in characteristic zero, $H''(r)\ne0$ generically; the marked
repeated root is then exactly double.  Birationality excludes a second
generic repeated root.
\end{proof}

Proposition~\ref{prop:universal-disc} is the boundary restriction of the
tangent-map core in Theorem~\ref{thm:tangent-core}: its critical equation is
$\gamma=0$.  The same theorem computes
\[
 \Omega_{\A^1_r/D_H}\simeq k[r]/(H''(r))\,dr,
 \qquad
 \operatorname{Fitt}_0\Omega_{\A^1_r/D_H}=(H''(r)),
\]
so differentiating the boundary restriction retains the full
scheme-theoretic Hessian decoration, including multiplicities.

\begin{lemma}[generic Morse slice]\label{lem:generic-morse-slice}
There is a finite set $\Sigma_H\subset\overline{k}$ such that
$H(W)-\sigma W$ is Morse for every
$\sigma\in\overline{k}\setminus\Sigma_H$: its $n-1$ critical points are
simple and have pairwise distinct critical values.
\end{lemma}

\begin{proof}
The critical points are the roots of $H'(W)-\sigma$.  They are simple unless
$\sigma$ belongs to the finite set
\[
 H'\bigl(V(H'')\bigr).
\]
If two distinct critical points $r_1,r_2$ have the same critical value, then
\[
 H'(r_1)=H'(r_2)=\sigma,\qquad
 r_1H'(r_1)-H(r_1)=r_2H'(r_2)-H(r_2).
\]
Thus the two points have the same image under the normalization map
\eqref{eq:disc-param}.  A finite birational morphism of curves is an
isomorphism away from a finite set, so only finitely many points of
$\A^1_r$ can occur in such a pair.  Excluding their finitely many
$s$-coordinates together with $H'(V(H''))$ proves the claim.
\end{proof}

\begin{theorem}[full symmetric monodromy]\label{thm:symmetric-monodromy}
The geometric Galois group of \eqref{eq:universal-pencil} over
$\overline{k}(s,t)$ and the arithmetic Galois group over $k(s,t)$ are both
$S_n$ in their natural action on the $n$ roots.
\end{theorem}

\begin{proof}
Choose $\sigma\notin\Sigma_H$ as in
Lemma~\ref{lem:generic-morse-slice}.  The classical monodromy theorem for a
Morse polynomial gives
\[
 \operatorname{Gal}\bigl(H(W)-\sigma W-T\,/\,\overline{k}(T)\bigr)=S_n
 \cite[Theorem~4.4.5]{serre-topics}.
\]
Indeed, the polynomial cover is connected, every finite branch point has
transposition inertia, and those inertia groups generate; a transitive group
generated by transpositions is $S_n$.

After $T=-t$, this cover is the restriction of
\eqref{eq:universal-pencil} to the vertical line $s=\sigma$.
Functoriality of geometric monodromy therefore includes its $S_n$ as a
subgroup of the two-parameter geometric group.  The latter is already a
subgroup of $S_n$, so equality holds.  The arithmetic group contains the
geometric group and is contained in $S_n$, hence is also $S_n$.
\end{proof}

For an admissible weighted map, on $D(C)$ put $s=BC$ and $t=cAC^2$.
Then $k(A,B,C)=k(s,t,C)$, a purely transcendental extension of $k(s,t)$.
Theorem~\ref{thm:symmetric-monodromy} therefore gives geometric and arithmetic
monodromy $S_n$ for the generic inverse cover of $G_H$.
